\documentclass[10pt,a4paper]{amsart}
\usepackage[margin=19mm]{geometry}
\usepackage{amsmath,amssymb,mathtools}
\usepackage[scr=rsfs,cal=euler]{mathalfa}
\usepackage{xfrac}
\usepackage[unicode,hidelinks,bookmarks=false]{hyperref}
\newcommand{\ie}{i.e.\ }
\newcommand{\cf}{cf.\ }
\newcommand{\resp}{respectively\ }
\newcommand{\Subsection}{\subsection}
\providecommand{\nobreakdash}{\nobreak\hbox{-}}
\setlength{\emergencystretch}{2em}
\multlinegap0pt
\usepackage{amssymb}
\usepackage{mathtools}
\usepackage{tikz}
\usepackage{cancel}
\usepackage[normalem]{ulem}
\newtheorem{lemma}{Lemma}
\newtheorem{proposition}{Proposition}
\usepackage{cleveref}
\crefname{lemma}{Lemma}{Lemmas}
\crefname{proposition}{Proposition}{Propositions}
\title{Geography and Deformations of $\mathbb{Z}_2^s$-Covers of General Type Over Weighted Projective Threefolds.}
\author{Patricio Gallardo, Jayan Mukherjee}
\date{Source: arXiv:2605.05665v1 (CC BY 4.0)}
\begin{document}
\maketitle

\noindent\emph{Source and licence.} Geography and Deformations of $\mathbb{Z}_2^s$-Covers of General Type Over Weighted Projective Threefolds.. Version arXiv:2605.05665v1 (CC BY 4.0), distributed by arXiv under the Creative Commons Attribution 4.0 International licence, http://creativecommons.org/licenses/by/4.0/, as recorded in the arXiv licence metadata for this exact version and shown on the abstract page https://arxiv.org/abs/2605.05665v1. Full licence notice: \texttt{licenses/arxiv-2605.05665-CC-BY-4.0.txt}.\par\smallskip

\noindent\textit{Editorial scope.} Counted unit: the article's proposition impossibility (source0.tex lines 5815--5838). The article's complete original proof of it is delivered with it (source0.tex lines 5840--6310, one \texttt{proof} environment of 471 lines). No mathematics is written, repaired or re-derived: the statement and the proof are the article's own lines, in the article's own notation, and no retained line is substituted or re-flowed.  Retained uncounted as the source-local context the proof needs: lines 3646--3665; lines 3871--3898.  Nothing was omitted from any retained span, so the package carries no omission record.  No retained line contains a citation, so the wrapper carries no bibliography and no bibliography entry.  No retained line contains a figure environment, an image-inclusion command or drawing code, so no image is delivered and the package declares no asset dependency.  Editorial formatting only, in addition: an AMS \texttt{amsart} preamble that restates the article's own declarations and macros used by the retained lines, namely the environments \texttt{lemma}, \texttt{proposition} on separate counters, together with the standard packages those lines need.  The retained environments print in this document as Lemma 1, Proposition 1 in order of appearance, whereas the article prints the same items with the numbering of its own full document.  The complete native source of the article as distributed in the arXiv e-print for this exact version is bundled unchanged as \texttt{sources/199975/source0.tex}.
\begin{lemma}\label{lemma: Fourier}


Let \( d: G \to \mathbb{Z}_{\ge 0} \) with \( d(0) = 0 \), and let \( l: G^* \setminus \{0\} \to \mathbb{Z}_{\ge 0} \). Define the Fourier transform of \( d \) as:
\[
\widehat{d}(\chi) = \sum_{x \in G} d(x) \, (-1)^{\chi \cdot x}.
\]

Suppose for all nontrivial \( \chi \neq 0 \):
\begin{equation}\label{l_chi}
\sum_{\substack{x \in G \\ \chi \cdot x = 1}} d(x) = 2 l(\chi). 
\end{equation}

Then for all \( \chi \neq 0 \):
\[
\widehat{d}(0) - \widehat{d}(\chi) = 4 l(\chi),
\]
where \( \widehat{d}(0) = \sum_{x \in G} d(x) \), and \( 0 \) denotes the trivial character.

\end{lemma}

\begin{lemma}[Fourier transform of convolution and the cubic moment identity]\label{lem:fourier-convolution-cubic}
Let $G=(\mathbb Z_2)^s$ written additively, and identify $G^*$ with $G$ via
$\chi(x)=(-1)^{\chi\cdot x}$, where $\chi\cdot x\in\mathbb Z_2$.
For a function $f:G\to\mathbb C$, define the (unnormalized) Fourier transform
\[
\widehat f(\chi)=\sum_{x\in G} f(x)\,(-1)^{\chi\cdot x},\qquad \chi\in G^*,
\]
and define convolution by
\[
(f*g)(z)=\sum_{\substack{x,y\in G\\ x+y=z}} f(x)g(y),\qquad z\in G.
\]
Then:
\begin{enumerate}
\item \textbf{(Convolution becomes pointwise product.)} For all $\chi\in G^*$,
\[
\widehat{f*g}(\chi)=\widehat f(\chi)\,\widehat g(\chi).
\]
\item \textbf{(Cubic moment at the origin.)} For any $d:G\to\mathbb C$,
\[
(d*d*d)(0)=\frac{1}{|G|}\sum_{\chi\in G^*}\widehat d(\chi)^3.
\]
In particular, if $d=\mathbf{1}_S$ is the indicator function of a subset $S\subseteq G$,
then
\[
(d*d*d)(0)=\#\{(x,y,z)\in S^3:\ x+y+z=0\}\ \in\ \mathbb Z_{\ge 0}.
\]
\end{enumerate}
\end{lemma}

\begin{proposition}\label{m = 1; k = 2; D = 12 impossibility}
Suppose $m = 1, k = 2$ and $D = 12$. Let $S:=\{g\in G : d(g)>0\}$ and $m_i = | g \in S, d_g = i|$.  Then there do not exist admissible solutions for 
\begin{enumerate}
    \item $s = 7, (m_1) = (12)$
    \item $s = 6, (m_1) = (12)$
    \item $s = 6, (m_1, m_2) = (10,1)$
    \item $s = 5, (m_1,m_3) = (9,1)$
    \item $s = 5, (m_1,m_2) = (8,2)$
    \item $s = 5, (m_1,m_2) = (10,1)$
    \item $s = 5, (m_1) = (12)$
   % \item $s = 4,  (m_1, m_2) = (5,2)$
    %\item $s = 4,  (m_1, m_3) = (6,1)$
    \item $s = 4, (m_1,m_2,m_3)=(7,1,1)$
    \item $s = 4, (m_1,m_3)=(9,1)$
    \item $s = 4, (m_2)=(6)$
    \item $s = 4, (m_1,m_2)=(2,5)$
    \item $s = 4, (m_1,m_2)=(4,4)$
    \item $s = 4, (m_1,m_2)=(6,3)$
    \item $s = 4, (m_1,m_2)=(8,2)$
    \item $s = 4, (m_1,m_4)=(8,1)$
    \item $s = 4, (m_1,m_2)=(10,1)$.
    
\end{enumerate}
\end{proposition}

\begin{proof}
\noindent\textit{Proof of $(1)$}
\begin{enumerate}
\item \textbf{First moment and the $\ell$--distribution.}
We have
\begin{equation}\label{eq:first-moment}
\sum_{\chi\in G}\ell_\chi = 2^{s-2}D = 384.
\end{equation}

Since $\ell_0=0$, we have
\[
\sum_{\chi\neq 0}\ell_\chi=384.
\]
There are $2^7-1=127$ nonzero $\chi$, and by assumption $\ell_\chi\ge 3$ for all
$\chi\neq 0$. Write
\[
A_\chi := \ell_\chi-3 \in \mathbb Z_{\ge 0}\qquad (\chi\neq 0).
\]
Then
\[
\sum_{\chi\neq 0} A_\chi
=
\sum_{\chi\neq 0}\ell_\chi - 3\cdot 127
=
384-381
=
3.
\]
In particular, the only possible
$\ell$--distributions are:
\[
(n_4,n_5,n_6)=(3,0,0),\quad (1,1,0),\quad (0,0,1),
\]
where $n_j:=\#\{\chi\neq 0:\ell_\chi=j\}$ and $n_3 = 124, 125, 126$ respectively.

\item \textbf{Fourier relation and the quadratic (Parseval) moment.}
Let $\widehat d:G\to \mathbb Z$ denote the (unnormalized) Fourier transform
\[
\widehat d(\chi):=\sum_{g\in G} d(g)(-1)^{\chi\cdot g}.
\]
By the Fourier relation in \Cref{lemma: Fourier}, for every $\chi\neq 0$,
\begin{equation}\label{eq:fourier-relation}
\widehat d(\chi)=\widehat d(0)-4\ell_\chi.
\end{equation}
Here $\widehat d(0)=\sum_{g} d(g)=D=12$, hence
\[
\widehat d(\chi)=12-4\ell_\chi
\qquad (\chi\neq 0).
\]
Therefore, when $\ell_\chi\in\{3,4,5,6\}$ we have
\[
\widehat d(\chi)\in\{0,-4,-8,-12\}.
\]

Now apply Plancherel/Parseval for the unnormalized Fourier transform:
\[
\sum_{\chi\in G}\widehat d(\chi)^2
=
|G|\sum_{g\in G} d(g)^2.
\]
Since $d(g)\in\{0,1\}$, we have $d(g)^2=d(g)$ and hence $\sum_g d(g)^2=D=12$.
As $|G|=2^7=128$, we obtain
\[
\sum_{\chi\in G}\widehat d(\chi)^2 = 128\cdot 12=1536.
\]
Separating the trivial character term $\widehat d(0)^2=12^2=144$ yields
\begin{equation}\label{eq:nontrivial-L2}
\sum_{\chi\neq 0}\widehat d(\chi)^2 = 1536-144=1392.
\end{equation}

On the other hand, using the possible $\ell$--distributions from Step~(1), we have
\[
\sum_{\chi\neq 0}\widehat d(\chi)^2
=
0\cdot n_3 + 16 n_4 + 64 n_5 + 144 n_6,
\]
so it must equal one of
\[
16\cdot 3=48,\qquad 16+64=80,\qquad 144.
\]
contradicting \eqref{eq:nontrivial-L2}.
\end{enumerate}

\noindent\textit{Proof of $(2)-(3)$}
We argue uniformly by the first and quadratic moments.

\begin{enumerate}
\item \textbf{First moment and the possible $\ell$--distributions.}
We have
\[
\sum_{\chi\in G}\ell_\chi = 2^{s-2}D = 192.
\]
Since $\ell_0=0$, we obtain $\sum_{\chi\neq 0}\ell_\chi=192$.
For $\chi\neq 0$, we have $2\ell_\chi=\sum_{\chi\cdot g=1} d(g)\le D=12$, hence
$\ell_\chi\le 6$. Together with $\ell_\chi\ge 3$, this implies
\[
\ell_\chi\in\{3,4,5,6\}\qquad(\chi\neq 0).
\]
Let
\[
n_j:=\#\{\chi\neq 0:\ell_\chi=j\},\qquad j=3,4,5,6.
\]
Then
\[
n_3+n_4+n_5+n_6=63,
\qquad
3n_3+4n_4+5n_5+6n_6=192.
\]
Subtracting $3\cdot 63=189$ from the second equation gives
\[
n_4+2n_5+3n_6=3.
\]
Thus the only possibilities are
\[
(n_4,n_5,n_6)=(3,0,0),\quad (1,1,0),\quad (0,0,1),
\]
with $n_3=60,61,62$ respectively.

\item \textbf{Fourier relation and the quadratic constraint from Plancherel.}
Let $\widehat d:G\to\mathbb Z$ denote the unnormalized Fourier transform
\[
\widehat d(\chi):=\sum_{g\in G} d(g)(-1)^{\chi\cdot g}.
\]
By the Fourier relation in \Cref{lemma: Fourier}, for every $\chi\neq 0$,
\[
\widehat d(\chi)=\widehat d(0)-4\ell_\chi.
\]
Since $\widehat d(0)=\sum_g d(g)=D=12$, we have
\[
\widehat d(\chi)=12-4\ell_\chi
\qquad (\chi\neq 0).
\]
Hence, for $\chi\neq 0$,
\[
\ell_\chi=3,4,5,6 \quad\Longrightarrow\quad
\widehat d(\chi)=0,-4,-8,-12,
\]
and therefore
\[
\widehat d(\chi)^2\in\{0,16,64,144\}.
\]
Using the possible $\ell$--distributions from Step~(1), we obtain
\[
\sum_{\chi\neq 0}\widehat d(\chi)^2
=
0\cdot n_3 + 16n_4 + 64n_5 + 144n_6
\in\{48,80,144\}.
\]

On the other hand, Plancherel's identity for the unnormalized Fourier transform gives
\[
\sum_{\chi\in G}\widehat d(\chi)^2
=
|G|\sum_{g\in G} d(g)^2,
\]
where $|G|=2^6=64$. Since $\widehat d(0)=D=12$, we may rewrite this as
\begin{equation}\label{eq:nontriv-energy-s6}
\sum_{\chi\neq 0}\widehat d(\chi)^2
=
64\sum_{g\in G} d(g)^2 - 12^2.
\end{equation}

We now compute $\sum_g d(g)^2$ in each of the two cases:

\begin{enumerate}
\item If $(m_1)=(12)$, then $\sum_g d(g)^2 = 12\cdot 1^2 = 12$,
so \eqref{eq:nontriv-energy-s6} gives $\sum_{\chi\neq 0}\widehat d(\chi)^2=64\cdot12-144=624$.
\item If $(m_1,m_2)=(10,1)$, then $\sum_g d(g)^2 = 10\cdot 1^2 + 1\cdot 2^2=14$,
so \eqref{eq:nontriv-energy-s6} gives $\sum_{\chi\neq 0}\widehat d(\chi)^2=64\cdot14-144=752$.
\end{enumerate}
In either case the value forced by Plancherel is $624$ or $752$, neither of which lies
in $\{48,80,144\}$. This contradicts the constraint coming from the possible
$\ell$--distributions.

\end{enumerate}



\noindent\textit{Proof of $(4)-(7)$} We argue uniformly by the first and quadratic moments.

\begin{enumerate}
\item[(i)] \textbf{First moment and the possible $\ell$--distributions.}
We have
\[
\sum_{\chi\in G}\ell_\chi = 2^{s-2}D = 96.
\]
For $\chi\neq 0$, we have $2\ell_\chi=\sum_{\chi\cdot g=1} d(g)\le D=12$, hence
$\ell_\chi\le 6$. Together with $\ell_\chi\ge 3$, this implies
\[
\ell_\chi\in\{3,4,5,6\}\qquad(\chi\neq 0).
\]
Let
\[
n_j:=\#\{\chi\neq 0:\ell_\chi=j\},\qquad j=3,4,5,6.
\]
Then
\[
n_3+n_4+n_5+n_6=31,
\qquad
3n_3+4n_4+5n_5+6n_6=96.
\]
Subtracting $3\cdot 31=93$ from the second equation gives
\[
n_4+2n_5+3n_6=3.
\]
Thus the only possibilities are
\[
(n_4,n_5,n_6)=(3,0,0),\quad (1,1,0),\quad (0,0,1),
\]
with $n_3=28,29,30$ respectively.

\item[(ii)] \textbf{Fourier relation and the quadratic constraint from Plancherel.}
Let $\widehat d:G\to\mathbb Z$ denote the unnormalized Fourier transform
\[
\widehat d(\chi):=\sum_{g\in G} d(g)(-1)^{\chi\cdot g}.
\]
By the Fourier relation in \Cref{lemma: Fourier}, for every $\chi\neq 0$,
\[
\widehat d(\chi)=\widehat d(0)-4\ell_\chi.
\]
Since $\widehat d(0)=\sum_g d(g)=D=12$, we have
\[
\widehat d(\chi)=12-4\ell_\chi
\qquad (\chi\neq 0).
\]
Hence, for $\chi\neq 0$,
\[
\ell_\chi=3,4,5,6 \quad\Longrightarrow\quad
\widehat d(\chi)=0,-4,-8,-12,
\]
and therefore
\[
\widehat d(\chi)^2\in\{0,16,64,144\}.
\]
Using the possible $\ell$--distributions from Step~(i), we obtain
\[
\sum_{\chi\neq 0}\widehat d(\chi)^2
=
0\cdot n_3 + 16n_4 + 64n_5 + 144n_6
\in\{48,80,144\}.
\]

On the other hand, Plancherel's identity for the unnormalized Fourier transform gives
\[
\sum_{\chi\in G}\widehat d(\chi)^2
=
|G|\sum_{g\in G} d(g)^2,
\]
where $|G|=2^5=32$. Since $\widehat d(0)=D=12$, we may rewrite this as
\begin{equation}\label{eq:nontriv-energy}
\sum_{\chi\neq 0}\widehat d(\chi)^2
=
32\sum_{g\in G} d(g)^2 - 12^2.
\end{equation}

We now compute $\sum_g d(g)^2$ in each of the four cases:
\begin{enumerate}
\item If $(m_1,m_2)=(8,2)$, then $\sum_g d(g)^2 = 8\cdot 1^2 + 2\cdot 2^2=16$,
so \eqref{eq:nontriv-energy} gives $\sum_{\chi\neq 0}\widehat d(\chi)^2=32\cdot16-144=368$.
\item If $(m_1,m_2)=(10,1)$, then $\sum_g d(g)^2 = 10\cdot 1^2 + 1\cdot 2^2=14$,
so \eqref{eq:nontriv-energy} gives $\sum_{\chi\neq 0}\widehat d(\chi)^2=32\cdot14-144=304$.
\item If $(m_1,m_2)=(12,0)$, then $\sum_g d(g)^2 = 12\cdot 1^2 = 12$,
so \eqref{eq:nontriv-energy} gives $\sum_{\chi\neq 0}\widehat d(\chi)^2=32\cdot12-144=240$.
\item If $(m_1,m_3)=(9,1)$, then $\sum_g d(g)^2 = 9\cdot 1^2 + 1\cdot 3^2=18$,
so \eqref{eq:nontriv-energy} gives $\sum_{\chi\neq 0}\widehat d(\chi)^2=32\cdot18-144=432$.
\end{enumerate}
In each case the value forced by Plancherel is one of
\[
368,\ 304,\ 240,\ 432,
\]
none of which lies in $\{48,80,144\}$. This contradicts the constraint coming from
the possible $\ell$--distributions.
\end{enumerate}


\noindent\textit{Proof of $(8)-(15)$}
We argue uniformly using the first, quadratic, and cubic moments.

\begin{enumerate}
\item[(i)] \textbf{First moment and the possible $\ell$--distributions.}
Since $s=4$ and $D=12$, the first moment identity gives
\[
\sum_{\chi\neq 0}\ell_\chi = 2^{s-2}D = 48.
\]
For $\chi\neq 0$, we have $2\ell_\chi=\sum_{\chi\cdot g=1} d(g)\le D=12$, hence
$\ell_\chi\le 6$. Together with $\ell_\chi\ge 3$, this implies
\[
\ell_\chi\in\{3,4,5,6\}\qquad(\chi\neq 0).
\]
Let
\[
n_j:=\#\{\chi\neq 0:\ell_\chi=j\},\qquad j=3,4,5,6.
\]
Since there are $15$ nontrivial characters, we have
\[
n_3+n_4+n_5+n_6=15,
\qquad
3n_3+4n_4+5n_5+6n_6=48.
\]
Subtracting $3\cdot 15=45$ from the second equation yields
\[
n_4+2n_5+3n_6=3.
\]
Thus the only possibilities are
\[
(n_4,n_5,n_6)=(3,0,0),\quad (1,1,0),\quad (0,0,1),
\]
with $n_3=12,13,14$ respectively.

\item[(ii)] \textbf{Quadratic moment (Plancherel).}
Let $\widehat d:G\to\mathbb Z$ denote the unnormalized Fourier transform
\[
\widehat d(\chi):=\sum_{g\in G} d(g)(-1)^{\chi\cdot g}.
\]
For $\chi\neq 0$, by \Cref{lemma: Fourier}, we have 
\[
\widehat d(\chi)=\widehat d(0)-4\ell_\chi.
\]
Since $\widehat d(0)=D=12$, it follows that
\[
\widehat d(\chi)=12-4\ell_\chi\qquad(\chi\neq 0),
\]
so that
\[
\ell_\chi=3,4,5,6
\quad\Longrightarrow\quad
\widehat d(\chi)=0,-4,-8,-12.
\]
Hence
\[
\sum_{\chi\neq 0}\widehat d(\chi)^2
=
16n_4+64n_5+144n_6
\in\{48,80,144\}.
\]

On the other hand, Plancherel's identity gives
\[
\sum_{\chi\in G}\widehat d(\chi)^2
=
|G|\sum_{g\in G} d(g)^2,
\qquad |G|=16.
\]
Since $\widehat d(0)=12$, we obtain
\begin{equation}\label{eq:s4-nontriv-energy}
\sum_{\chi\neq 0}\widehat d(\chi)^2
=
16\sum_{g\in G} d(g)^2 - 12^2.
\end{equation}
A direct computation shows that for the cases
\[
(m_1,m_2,m_3) = (7,1,1),\ (m_2) = (6),\ (m_1,m_2) = (2,5),\ (4,4),\ (m_1, m_2) = (8,2),\ (m_1, m_4) = (8,1)
\]

the value forced by \eqref{eq:s4-nontriv-energy} does not lie in
$\{48,80,144\}$, yielding an immediate contradiction. Thus these cases are
ruled out by the quadratic moment.

\item[(iii)] \textbf{Cubic moment.}
It remains to consider the cases $(m_1,m_3)=(9,1)$ and $(m_1,m_2)=(6,3)$ and $(m_1, m_2)) = (10, 1)$. For the first two, \eqref{eq:s4-nontriv-energy} gives
\[
\sum_{\chi\neq 0}\widehat d(\chi)^2=144.
\]
Hence the $\ell$--distribution must be $(n_4,n_5,n_6)=(0,0,1)$.
For this distribution we have
\[
\sum_{\chi\neq 0}\widehat d(\chi)^3 = (-12)^3=-1728,
\]
and therefore
\[
\sum_{\chi\in G}\widehat d(\chi)^3
=
12^3-1728=0.
\]
Using the identity from \Cref{lem:fourier-convolution-cubic},
\[
\sum_{\chi\in G}\widehat d(\chi)^3
=
|G|\,(d*d*d)(0),
\]
we conclude that $(d*d*d)(0)=0$. However, in both cases $(9,1)$ and $(6,3)$ the support $S=\{g:d(g)\ge 1\}$ has cardinality $|S|>8$. In $(\mathbb Z_2)^4$, any subset of size $>8$ contains elements $a,b$ with $a+b\in S$, and hence contributes positively to
$(d*d*d)(0)$. This contradiction shows that these two cases are impossible as
well.

\item[(iv)] \textbf{Fractional values using inverse Fourier Transform}

\medskip
We finally deal with the case $(m_1,m_2)=(10,1)$.
In this case we have
\[
\sum_{g\in G} d(g)^2 = 10\cdot 1^2 + 1\cdot 2^2 = 14.
\]
By Plancherel's identity for the unnormalized Fourier transform on $G$ with
$|G|=16$, we obtain
\[
\sum_{\chi\in G}\widehat d(\chi)^2
=
16\sum_{g\in G} d(g)^2
=
16\cdot 14
=
224.
\]
Since $\widehat d(0)=\sum_g d(g)=D=12$, it follows that
\[
\sum_{\chi\neq 0}\widehat d(\chi)^2
=
224-12^2
=
80.
\]

From the classification of $\ell$--distributions established earlier for
$D=12$, $s=4$, and $\ell_\chi\ge 3$, the three possible types give
\[
\sum_{\chi\neq 0}\widehat d(\chi)^2\in\{48,80,144\}.
\]
Hence the value $80$ forces the $\ell$--distribution to be of \emph{Type~B}.
Equivalently, there exist distinct nontrivial characters $u\neq v$ such that
\[
\widehat d(u)=-4,\qquad
\widehat d(v)=-8,\qquad
\widehat d(\chi)=0\ \text{for all other }\chi\neq 0,
\]
and $\widehat d(0)=12$.

We now apply the inverse Fourier transform. For every $x\in G$,
\[
d(x)
=
\frac{1}{16}\sum_{\chi\in G}\widehat d(\chi)(-1)^{\chi\cdot x}
=
\frac{1}{16}\Bigl(12-4(-1)^{u\cdot x}-8(-1)^{v\cdot x}\Bigr).
\]

We evaluate this expression according to the values of
$(u\cdot x,v\cdot x)\in(\mathbb Z_2)^2$.

\begin{enumerate}
\item If $u\cdot x=0$ and $v\cdot x=0$, then
\[
d(x)=\frac{12-4-8}{16}=0.
\]

\item If $u\cdot x=1$ and $v\cdot x=0$, then
\[
d(x)=\frac{12+4-8}{16}=\frac12.
\]

\item If $u\cdot x=0$ and $v\cdot x=1$, then
\[
d(x)=\frac{12-4+8}{16}=1.
\]

\item If $u\cdot x=1$ and $v\cdot x=1$, then
\[
d(x)=\frac{12+4+8}{16}=\frac32.
\]
\end{enumerate}

In particular, $d(x)$ takes the nonintegral values $\frac12$ and $\frac32$,
contradicting the assumption that $d(x)\in\mathbb Z_{\ge 0}$ for all $x\in G$.
Therefore the case $(m_1,m_2)=(10,1)$ cannot occur.



\end{enumerate}

%The only remaining possibilities for $D=12$, $s=4$, and $\ell_\chi\ge 3$ are

\end{proof}

\end{document}
